% Page budget: 1.55 pages | Owns (README ownership table): augmented model, encoder, closed-loop criterion, added-state initialisation (PS2), optimiser and checkpoint selection; models M_U, M_PS2.
% WRITING GUIDE
% Purpose: Define the final augmentation architecture and explain how it is trained inside the known feedback loop.
% Start here: Current implementation decisions D-129, D-140, D-141, D-180, D-220, D-233, and D-237 in docs/decisions.md.
% Supporting material: Quinten's 18 September outline feedback, the Meeting-audit synthesis, Thesis-writeup/Documentation/TRAINING-DESIGN.md, and docs/writeup figures.
% Mandatory papers: Read Hoekstra's 2025 LFR augmentation paper, the 2026 first-principles augmentation paper, the encoder-initialization paper, and Kessels' Chapter 5 before writing this section.
% Research-plan anchor: Preserve Aspect 2's dynamic parallel augmentation objective; document the departure from open-loop training and earlier state routing.
% Current decisions: Separate physical and additional states, keep the controller outside the model, score the complete training window without burn-in, and use one network for the physical correction and added-state update.
% Choices to justify: Final reported routing, added-state dimension, ANN width and depth, zero initialization, encoder history, closed-loop objective, rollout length, full-window scoring, and validation horizon. Use the configuration of the reported thesis run, not a superseded routing experiment.
% Horizon check: Treat encoder history, scored training window, joint-estimation horizon, and free-run validation horizon separately; relate candidates to relevant stable modes and test sensitivity to slower dynamics and free-integrator accumulation.
% Implementation audit: Read scripts/gantry/gantry_interconnect_dynamic.py, gantry_dynamic/{model,config,controller,data,training}.py, and model_augmentation/fit_systems/{interconnect,closed_loop,pre_encoder}.py.
% Reference check: Verify sources for parallel LFR augmentation, encoder initialization, closed-loop state extension, and any claimed architectural precedent against the original paper or thesis.
% Cautions: Earlier routing plans are superseded; distinguish the truth-only hidden absorber from learned states; closed-loop fit alone does not prove autonomous plant recovery.
% Figures: Absorber schematic, author's updated architecture, and thesis-owned common-reference loops. No residual panel or horizon figure.
% Section is finished when: Every signal has a defined frame and timing, the RK4 boundary is explicit, and the described architecture matches the final code.
%
% SOURCE MAP
% Hoekstra et al. 2025 EJC: dynamic-parallel topology, additional states, encoder-based truncated simulation.
% Hoekstra et al. 2026 Automatica submission: extended LFR formulation, baseline-preserving model/encoder initialization, joint identification context.
% Hoekstra et al. 2026 encoder paper: reconstructability-based W^b initialization; it does not identify W^a or test dynamic augmentation.
% Kessels 2025 thesis: direct closed-loop rollout, controller equations, and reconstruction of the machine controller state for each window.
% Drenth 2025 thesis: LPV-specific augmentation precedent; not the unrestricted MLP structure used here.
% This project: gantry specialization, residual-loop implementation, exact routing, RK4 boundary, and verification.
%
% MUST ESTABLISH (agreed with Dirk, 2026-10-06; step 2a of the thesis-section skill)
% Contribution delivered: the gantry realisation of S-DP augmentation (CT self-scheduled LPV-LFR baseline
%   inside a DT parallel model, joint estimation in the ten combinations) and its identification under
%   feedback. Not new: the S-DP topology (Hoekstra), augmentation of self-scheduled LPV-LFR with added
%   states (Drenth Eq. 5.2), closed-loop training with a known controller (Kessels, inspiration only).
% III-A Augmented model
%  1. Dynamic, not static: static augmentation adds no states (EJC Sec. 2, "flexible modes"); refitting
%     theta_base adds none either. Argue from the general case, no absorber.
%  1b. The augmented model is Hoekstra's LFR augmentation structure (EJC Eq. 4, 2026 Eq. 6) with a FIXED
%     interconnection (EJC Sec. 3.2) realising S-DP; its baseline block is Section II's self-scheduled
%     LFR with Delta(Y) = Y I6: an LFR inside an LFR, which is what the section title means.
%     Well-posedness: acyclic graph (EJC Thm. 1, 2026 Cond. 6) plus Section II's M(Y) condition.
%  2. Parallel, not series. LEAD REASON (Dirk, 2026-10-06): the orthogonal-by-construction method that
%     Section IV builds on is defined for the additive (parallel) structure (Gyorok 2026 abstract, Sec. 1;
%     D-003), so parallel lets us use that existing work. Supporting: Hoekstra gives no universal rule,
%     the structure follows prior knowledge (2026 Sec. 1, 6.3); the baseline stays the explicit nominal
%     term (research plan Aspect 2); parallel needs only a feedforward network for baseline-preserving
%     init, series needs ResNets (EJC Sec. 4.3, 2026 Sec. 5.4, 6.3).
%  3. State augmentation only, h_aug = 0: output P^T q is exact kinematics (cf. Kessels Remark 5.3). TODO reason.
%  4. Correction on all physical rows and added states: S-DP definition (EJC), X and Y reachable (D-103);
%     marginal-mode risk on free axes stated. TODO: why correct position rows rather than keep kinematics.
%     Context: Hoekstra 2026 Sec. 7 (F1Tenth) detaches the free integrators and learns only the
%     input-to-velocity map because the full state is measured; here only positions are measured.
%     One network for f_aug and g_aug is an implementation statement (D-180), no scientific claim.
%  5. DT correction after the CT RK4 step (input held, M(Y) per stage, affine normalisation). TODO reason
%     versus force injection in the CT derivative.
%  6. Unrestricted nonlinear augmentation (depends on Y through the state) versus LPV-structured
%     augmentation (Drenth): departure from the research plan's "in LPV-LFR form". TODO reason. Acyclic
%     routing adds no instantaneous loops; baseline well-posedness is Section II's.
%  7. Zero output layer reproduces the baseline from the same physical initial state (Hoekstra 2026
%     Eq. 29, D-237); the split is not unique (Sec. 5.2), handed to Section IV. Added states are memory
%     and model order, defined up to a state transformation, no physical coordinates.
% III-B Encoder
%  8. Encoder needed: only positions measured; truncated windows for cost and stability (Beintema).
%  9. W^b from the reconstructability map of a ZOH linearisation at an equilibrium (encoder paper
%     Eqs. 15-17, 30-32), not fitted (Hoekstra 2026 Eq. 30); W^a random. Claim only better starting
%     values and faster convergence, final accuracy comparable (encoder paper Sec. 4.4). The paper
%     assumes a stable baseline; ours is marginally stable (free X, Y), the map exists because the
%     positions make the linearisation observable. The paper treats co-estimating theta_base as out of
%     scope (footnote 1); we co-estimate. Precedent for linearising at zero: encoder paper Sec. 4.2
%     ("around the 0 point"), no reason given. TODOs: why Y = 0; disclose nominal parameters in W^b
%     versus detuned start; why model-based rather than data-based init (candidate: cost, Sec. 4.4);
%     Kaiming versus Xavier (todo only).
% 10. Encoder estimates the full state incl. measured positions (SUBNET). TODO: mention Kessels Remark 5.3?
% 11. Keep encoder history, scored window and validation horizon apart; values in Setup. The
%     joint-estimation sensitivity horizon belongs to Section IV.
% III-C Closed-loop identification
% 12. Known-controller, reference-driven output-error identification, inspired by Kessels (Eqs. 5.12,
%     5.13; Remark 5.6 open-loop attempt failed). Related to the indirect approach (Forssell and Ljung
%     Sec. 5.2) as a linear-theory motivation, not a guarantee; direct identification is not called
%     inherently biased. Free X and Y axes: the SUBNET consistency argument that Hoekstra inherits
%     (2026 Sec. 5.5) assumes an incrementally exponentially output stable data-generating system
%     (Beintema 2023 Condition 1); the open-loop gantry is not (free integrators), the loop with the
%     stabilising controller can be. Precondition, not a proof, for our nonlinear model. Same principle
%     as Kessels (closed-loop simulation of an unstable open-loop system), scoped to marginal stability.
% 13. Our formulation (ours, "we"): u_hat = u + C x^c + D(y - y_hat). Derivation by subtracting the
%     machine's controller equations, equal to the shared-r, shared-u_ff loop of the figure under a
%     compatible controller, signals and initial state (assumption; whether the data meet it is a Setup
%     or Results check: 4 kHz controller on recorded r - y versus recorded feedback force).
%     Main consequence: x^c_tau = 0 per window, no controller initialisation; consistent with the
%     encoder (model state and model controller both start on the machine's condition at tau). Contrast:
%     Kessels' direct form needs the machine controller state (Remark 5.4). Error before tau is not
%     carried in; the free-run validation does carry it.
% 14. Step order from no plant feedthrough: no algebraic loop, no added delay.
% 15. Every window sample scored (as Eq. 22); theta_base, theta_aug, eta joint, controller fixed.
%     Joint estimation framing (Dirk, 2026-10-06): the goal is to estimate theta_base jointly, in the ten
%     combinations. Obstacle from Jan's work: the split is not unique (2026 Sec. 5.2) and approximately
%     initialised parameters stay near their start (EJC Sec. 4.3, 2026 Sec. 6.3). Handover: Section IV
%     extends Jan's method with orthogonality so that the estimate of theta_base is UNIQUE (this
%     strength, not "recovers the true parameters"). State what interpretable means: the estimated
%     parameters keep their physical meaning; cite Gyorok 2026 Sec. 3 ("With non-unique theta_b, the
%     baseline model could lose its physical meaning"). Operational form only; the general definition
%     is the Introduction's (todo there). "Negation" and Kessels' definition (Sec. 5.3.1.3) belong to
%     Section IV; his Sec. 6.2 recommendation to the Introduction gap. Do NOT mention Jan's parameter regularisation anywhere.
%     Do NOT say explicitly that Section III's formulation is the unconstrained comparison arm; let
%     Section IV extend it naturally.
% Setup pattern for hyperparameters (2026 Sec. 7): honest basis per value (physical insight and
%     trial-and-error, line-search, inherited from earlier results).
% 16. Closed-loop fit can mask plant error; plant-domain checks and controller transfer in Results.
% Notation: u = P u_act as in Sections II and IV; no "(eom)" shorthand; controller matrices distinct from A_c(Y).
% III-D Added-state initialisation (PROPOSED, cycle 07, inferred from the README ownership row III and
%   source-map row III-D; no agreed block exists)
% 17. Problem: at the zero start the added states do not reach the output, so V gives W^a and the added-state
%     rows of phi_aug no gradient; plain training afterwards learns static corrections or real lags (evidence
%     Section VI; PS2-METHOD Sec. 2.2: the gate alone does not argue the need).
% 18. PS2 (ours) adapts two-stage regression (Hefny Sec. 3) and 2SR-then-BPTT (Downey Sec. 4.2): S1 fits the
%     encoder's added-state rows and a temporary decoder to the current model's closed-loop residual over the
%     encoder horizon; S2 fits g_aug to propagate the encoder states one sample; S3 = Section III-C unchanged.
% 19. Screen and stopping rules as PS2Spec states them, values in Section V; M_PS2 named in its own paragraph.
% Restrictions (README PS2 sketch): about two displays and six sentences plus one per rule; no optimiser
%     schedule, calibration bookkeeping beyond naming the compared quantity, or gradient mechanics.
%
% CORRECTIONS AND HAND-OVERS (cycle 07)
% - Resolved: W^a is Xavier uniform, gain 1, in the thesis runs (D-239, _THESIS_FIXED wa_encoder_init); the
%   earlier Kaiming todo is removed. Resolved: refs.bib now holds gyorok2026obc, forssell1999revisited,
%   hefny2015supervised, downey2017psrnn (with CHECK comments), so the missing-bib todo is removed.
%   Resolved: model-based versus data-based W^b (reason from DOC/AUGMENTATION-INITIALIZATION.md Sec. 5.1).
% - Block item 13 corrected: x^c_tau = 0 is a definition of the residual-form initial condition
%   (FS/closed_loop.py), not Kessels' Remark 5.4 (README open conflicts, resolved: code wins).
% - Block item 9 corrected: the encoder history includes the window start (na_right = 1); Hoekstra 2026 Eq. (22e)
%   is past-only, the encoder paper's reconstructability map Eqs. (11) to (17) includes the current sample.
% - Notation: Gamma_n for the input-to-state map (r_n clashes with the reference r_k); kappa scales the S1 target.
% - Notation: theta_enc (not eta) for the encoder, as Section IV uses; the encoder histories are bold y_tau,
%   u_tau (xi is the training coordinate of Section IV); the machine controller state is x^fb (s clashes with
%   the scale s of eq:mdelta_map); the Toeplitz matrix is H_n (T_n clashes with T_x, T_u, T_y); the decoder
%   is lambda (h clashes with h_base, d with the offset d).
% - For Section V (05_setup.tex, not edited): the hyperparameter table needs the PS2 rows (S1 cap and minimum,
%   check cadence, plateau tolerances, screen threshold bar-gamma and its update, S2 minimum and cap, decoder
%   width); the bullet "Free-run validation (referenced from Section III-C)" is now defined in Section III-C,
%   Section V keeps only its records; the controller-rate bullet should carry the measured loop difference
%   (todo in Section III-C). 99_appendix.tex app:lfr-wellposed is consistent with eq:mdelta_map (unchanged).
\section{Dynamic LPV-LFR Augmentation}\label{sec:augmentation}
\ifdraft
\begin{itemize}
\item The section extends the baseline with learned dynamics and identifies the result from closed-loop records; it names its four subsections (style profile)
\item Contribution share: the realisation on the self-scheduled gantry baseline with joint estimation of $\theta_{\mathrm{base}}$, its identification in closed loop with the known controller, and the PS2 initialisation (MUST ESTABLISH; Introduction wording (?))
\end{itemize}
\fi

Given the baseline \eqref{eq:delta}, \eqref{eq:lfr_G}, our goal is a model
that also captures the dynamics the baseline omits, identified from records
measured under feedback, with $\theta_{\mathrm{base}}$ estimated jointly in the
combinations of Section~\ref{sec:params}.
To this end, Section~\ref{sec:aug_structure} defines the augmented model,
Section~\ref{sec:aug_encoder} the encoder that estimates its initial state,
Section~\ref{sec:aug_closed_loop} the closed-loop identification problem, and
Section~\ref{sec:aug_ps2} the initialisation of the added states.
Together they extend dynamic parallel LFR augmentation to a closed-loop,
self-scheduled state-space setting with learned additional states, the second
contribution of Section~I.
The values of the symbols defined here are given in Section~\ref{sec:setup}.
\todo{Missing: contribution wording. This section also delivers the joint estimation of $\theta_{\mathrm{base}}$ in the ten combinations and the PS2 initialisation, which the second contribution does not name. Candidate: ``\dots with learned additional states, initialised from data, and with joint estimation of the identifiable physical combinations'' (README ownership row III; D-236, D-240).}
\todo{VERIFY keys: \texttt{hoekstra2025lfr} = EJC 86 (2025) 101304,
\texttt{hoekstra2026lfr} = arXiv:2602.17297 (submitted to Automatica).
refs.bib and docs/references.md rows 46 and 130 now agree.}

\subsection{Augmented model}\label{sec:aug_structure}
\ifdraft
\begin{itemize}
\item If the omitted effects are dynamic, the six baseline states cannot represent them; refitting $\theta_{\mathrm{base}}$ or a static augmentation adds no states (EJC Sec.~2)
\item The model is Hoekstra's LFR augmentation structure with a fixed interconnection realising S-DP; its baseline block is the LFR of Section~II (EJC Sec.~3.2, 2026 Table~1)
\item Parallel because the orthogonal-by-construction method of Section~IV is defined for the additive structure (Gy\"or\"ok 2026 Abstract)
\item The interconnection is acyclic, so it adds no algebraic loop (EJC Thm.~1)
\item $f_{\mathrm{base}}$ is one RK4 step with the input held and $M(Y)$ per stage, so it stays self-scheduled
\item Normalisation displayed: affine, statistics of the training records, state statistics from $P^{-\top}y$ and its differences, unlike Hoekstra Eq.~(28); $\phi_{\mathrm{aug}}$ acts on normalised signals and supplies $f_{\mathrm{aug}}$ and $g_{\mathrm{aug}}$
\item Realisation choices in one paragraph: plain feedforward network (EJC Sec.~4.3); correction on all physical rows, $X$ and $Y$ included (?: why the position rows); correction after the RK4 step (?); no LPV structure, unlike Drenth Eq.~(5.2) (?); output map not learned (?)
\item A zero output layer reproduces the baseline from the same physical initial state (2026 Eq.~(29))
\item Added states have no physical coordinates; the baseline/network split is not unique (2026 Secs.~5.2, 6.3); interpretable means the estimated $\theta_{\mathrm{base}}$ keeps its physical meaning (Gy\"or\"ok Sec.~3); Section~IV makes the estimate unique
\end{itemize}
\fi

If the effects that the baseline omits are dynamic, its six states cannot
represent them.
Refitting $\theta_{\mathrm{base}}$ adds no states.
A static augmentation, which corrects the state update without extending the
state, adds none either~\cite[Sec.~2]{hoekstra2025lfr}.
The augmentation therefore needs states of its own.

We adopt the LFR-based augmentation structure of Hoekstra et al., in which a
fixed interconnection connects a baseline block and a learning
block~\cite[Sec.~3.2]{hoekstra2025lfr}.
With it, we realise the dynamic-parallel structure
of~\cite[Table~1]{hoekstra2026lfr}.
We choose the parallel structure because the orthogonal-by-construction
parametrisation that Section~\ref{sec:obc} extends is defined for additive
augmentation~\cite[Abstract]{gyorok2026obc}.
The baseline block is the self-scheduled LFR of Section~\ref{sec:lfr}, so the
augmented model is an LFR whose baseline block is itself an LFR in $Y$:
\begin{equation}
\begin{aligned}
 \tilde x_{k+1}
 &=f_{\mathrm{base}}(\theta_{\mathrm{base}},\tilde x_k,u_k)
   +f_{\mathrm{aug}}(\theta_{\mathrm{aug}},\hat x_k,u_k),\\
 \bar x_{k+1}
 &=g_{\mathrm{aug}}(\theta_{\mathrm{aug}},\hat x_k,u_k),\\
 \hat y_k&=h_{\mathrm{base}}(\tilde x_k),
\end{aligned}
\label{eq:aug_transition}
\end{equation}
where $\tilde x_k=\col(q,\dot q)\in\R^6$ is the physical state at sample $k$,
$\bar x_k\in\R^{n_{\bar x}}$ the added state,
$\hat x_k=\col(\tilde x_k,\bar x_k)$, $u_k=Pu_{\mathrm{act},k}\in\R^3$ the
generalised force of \eqref{eq:Ptransform}, $\hat y_k\in\R^3$ the predicted
actuator positions, $f_{\mathrm{base}}:\R^{10}\times\R^6\times\R^3\to\R^6$ the
baseline transition, $\theta_{\mathrm{aug}}$ the parameters of the learned maps
$f_{\mathrm{aug}}:\R^{6+n_{\bar x}}\times\R^3\to\R^6$ and
$g_{\mathrm{aug}}:\R^{6+n_{\bar x}}\times\R^3\to\R^{n_{\bar x}}$, and
$h_{\mathrm{base}}(\tilde x)=P^\top q$ the kinematic output map.
No learned map feeds the baseline block or itself within a step, so the
interconnection adds no algebraic loop~\cite[Thm.~1]{hoekstra2025lfr}.

The baseline transition $f_{\mathrm{base}}$ is one fourth-order Runge-Kutta
(RK4) step of \eqref{eq:eom} over the sample period $T_s$, with the input held
over the step.
Each stage evaluates $M(Y)$ at its own state, so $f_{\mathrm{base}}$ remains
self-scheduled.
The model runs in coordinates normalised with the statistics of the training
records.
Unlike Hoekstra et al., who only scale the signals and take the state scale
from a simulation of the baseline~\cite[Sec.~5.3]{hoekstra2026lfr}, we centre
and scale.
The state statistics come from the recorded positions $P^{-\top}y$ and their
first differences, because only positions are measured.
In these coordinates, the maps of \eqref{eq:aug_transition} act as
\begin{equation}
\begin{aligned}
 &\tilde x^{\mathrm n}=T_x(\tilde x-\mu_x),\quad
  u^{\mathrm n}=T_u(P^{-1}u-\mu_u),\\
 &y^{\mathrm n}=T_y(y-\mu_y),\\
 &f^{\mathrm n}_{\mathrm{base}}(\theta_{\mathrm{base}},\tilde x^{\mathrm n},u^{\mathrm n})\\
 &\quad=T_x\bigl(f_{\mathrm{base}}(\theta_{\mathrm{base}},T_x^{-1}\tilde x^{\mathrm n}+\mu_x,u)-\mu_x\bigr),\\
 &\begin{bmatrix}T_xf_{\mathrm{aug}}(\theta_{\mathrm{aug}},\hat x,u)\\
   g_{\mathrm{aug}}(\theta_{\mathrm{aug}},\hat x,u)\end{bmatrix}
  =\phi_{\mathrm{aug}}\bigl(\theta_{\mathrm{aug}},\col(\tilde x^{\mathrm n},\bar x),u^{\mathrm n}\bigr),
\end{aligned}
\label{eq:aug_norm}
\end{equation}
where $y\in\R^3$ is the recorded output $\qact$, $\mu_x\in\R^6$ and
$\mu_u,\mu_y\in\R^3$ are channel means, $T_x$, $T_u$ and $T_y$ are diagonal
matrices of inverse channel standard deviations, and
$\phi_{\mathrm{aug}}:\R^{6+n_{\bar x}}\times\R^3\to\R^{6+n_{\bar x}}$ is a
multilayer perceptron with tanh hidden layers and a linear output layer.

The learned part rests on five realisation choices.
A plain feedforward network suffices, because the baseline is augmented
additively~\cite[Sec.~4.3]{hoekstra2025lfr}.
As the dynamic-parallel structure prescribes, $f_{\mathrm{aug}}$ corrects every
physical row~\cite[Sec.~2]{hoekstra2025lfr}, $X$ and $Y$ included.
On these axes a correction meets no stiffness in
\eqref{eq:damping_stiffness_matrices}, so a persistent correction accumulates in
position.
The correction is added after the RK4 step, so $f_{\mathrm{aug}}$ is a
discrete-time correction at the sample rate.
The network depends on $Y$ through $\hat x$ but has no LPV structure, unlike
the augmented LPV-LFR of Drenth~\cite[Eq.~(5.2)]{drenth2025thesis}.
The output map is not augmented.
\todo{Missing: why the position rows are corrected rather than kept as the integrals of the velocities. Context: Hoekstra et al.\ detach the free integrators of their vehicle and learn only the input-to-velocity map, which needs measured velocities~\cite[Secs.~7.2, 7.4]{hoekstra2026lfr}; here only positions are measured. D-103 requires only that $X$ and $Y$ are reachable.}
\todo{Missing: reason for correcting the state after the RK4 step rather than adding a learned force inside the continuous-time derivative. No decision entry gives a reason other than gradient magnitude.}
\todo{Missing: reason for an unrestricted network rather than an LPV-structured augmentation. The research plan (Aspect~2) proposed the interconnection in LPV-LFR form; this is a departure, and D-017 still requires an LPV-structured augmentation (README open conflicts).}
\todo{Missing: reason for state augmentation only ($h_{\mathrm{aug}}=0$). Candidate: $h_{\mathrm{base}}$ is the exact kinematic map of the measured positions, the premise under which Kessels initialises positions from the output~\cite[Remark~5.3]{kessels2025ai}; an output augmentation depending on the input would also close an algebraic loop with the controller~\cite[Remark~5.1]{kessels2025ai}.}

\begin{figure}[t]
 \centering
 \resizebox{\columnwidth}{!}{\input{figures/tex/augmentation_structure}}
 \caption{Augmented model. The physics step $f_{\mathrm{base}}$ and the
 network $\phi_{\mathrm{aug}}$ act in parallel on $\hat x_k$;
 $\phi_{\mathrm{aug}}$ supplies $f_{\mathrm{aug}}$ and $g_{\mathrm{aug}}$,
 and the output map is not learned ($h_{\mathrm{aug}}=0$). The encoder
 $\psi$ sets $\hat x$ at each window start. During training, $u_k$ is the
 closed-loop input of Section~\ref{sec:aug_closed_loop}.
 \todo{No setpoint generator currently for the inputs. to cramped together. We could make this a two column figure.}
 }
 \label{fig:aug_structure}
\end{figure}

Training starts from the baseline (Fig.~\ref{fig:aug_structure}).
The output layer of $\phi_{\mathrm{aug}}$ is initialised at zero,
\begin{equation}
 W_L=0,\quad b_L=0
 \quad\Longrightarrow\quad
 f_{\mathrm{aug}}=0,\quad g_{\mathrm{aug}}=0,
 \label{eq:aug_zero_init}
\end{equation}
where $W_L$ and $b_L$ are the weight matrix and bias vector of that layer.
From the same physical initial state, the augmented model then reproduces the
baseline, the initialisation condition of~\cite[Eq.~(29)]{hoekstra2026lfr}.
\todo{Missing: attribution of the all-zero output layer. It follows the authors' public dynamic-parallel code (D-237); \cite[Sec.~5.4.3]{hoekstra2026lfr} also allows a random linear part. Ask Jan which produced the reported S-DP results.}

The parameters of the augmented model are not unique.
An invertible transformation of $\bar x$, absorbed by $\phi_{\mathrm{aug}}$,
leaves the input-output behaviour unchanged, so the added states have no
physical coordinates.
When $\theta_{\mathrm{base}}$ is estimated jointly with $\theta_{\mathrm{aug}}$
(Section~\ref{sec:aug_closed_loop}), $f_{\mathrm{aug}}$ can also learn changes
that a different $\theta_{\mathrm{base}}$ would
produce~\cite[Sec.~5.2]{hoekstra2026lfr}.
The split between baseline and network is then not unique.
In the experiments of Hoekstra et al., approximately initialised physical
parameters stay close to their start, while the network learns dynamics that
the baseline could represent~\cite[Sec.~6.3]{hoekstra2026lfr}.
With a non-unique $\theta_{\mathrm{base}}$, the baseline can lose its physical
meaning~\cite[Sec.~3]{gyorok2026obc}.
Here, the estimated combinations are interpretable if they keep that meaning,
which requires a unique estimate.
Section~\ref{sec:obc} extends the method so that the estimate is unique.

\subsection{Encoder and initial state}\label{sec:aug_encoder}
\ifdraft
\begin{itemize}
\item Each window starts from a state of which only the positions are measured; an encoder estimates it from the input-output history (Beintema)
\item Windows starting from an estimated state keep the cost independent of the record length and stabilise the optimisation (Beintema Secs.~2, 3)
\item The encoder is a linear map on the scaled histories plus a nonlinear correction on the normalised ones (encoder paper Eq.~(8)); the history includes the window start, unlike 2026 Eq.~(22e)
\item The encoder estimates all states, the measured positions included (?: Kessels Remark~5.3)
\item $W^b$ is the reconstructability map of the baseline linearised at rest, $Y=0$, nominal parameters, ZOH (encoder paper Eqs.~(15) to (17)); computed, not fitted as in 2026 Eq.~(30), because the linearisation is known
\item The positions make the linearisation observable although the baseline is marginally stable
\item $W^a$ Xavier as in 2026 Eq.~(31); trained jointly; better starting values and convergence speed, not final accuracy, shown for a stable, not co-estimated baseline (encoder paper Secs.~2, 4.4)
\item Open: why $Y=0$; nominal parameters in $W^b$ under a detuned start (?)
\end{itemize}
\fi

Each training window starts from a state of which only the positions are
measured.
Following Beintema et al.~\cite{beintema2023subnet}, an encoder estimates this
state from the recorded input and output history.
Simulating short windows from encoded states keeps the cost of a gradient step
independent of the record length and makes the optimisation more
stable~\cite[Secs.~2, 3]{beintema2023subnet}.
The encoder is a linear map with a nonlinear
correction~\cite[Eq.~(8)]{hoekstra2026encoder}.
Its history includes the window start $\tau$, as in the reconstructability map
below~\cite[Eqs.~(11) to (17)]{hoekstra2026encoder}.
The encoder of~\cite[Eq.~(22e)]{hoekstra2026lfr} reads past samples only.
The linear part acts on scaled histories and the correction on normalised ones:
\begin{equation}
\begin{aligned}
 \begin{bmatrix}T_x(\tilde x_{\tau|\tau}-\mu_x)\\ \bar x_{\tau|\tau}\end{bmatrix}
 &=\psi(\theta_{\mathrm{enc}},\mathbf y_\tau,\mathbf u_\tau)\\
 &=\begin{bmatrix}W^b\\ W^a\end{bmatrix}
   \begin{bmatrix}\bar T_y\mathbf y_\tau\\ \bar T_u\mathbf u_\tau\end{bmatrix}
  -\begin{bmatrix}T_x\mu_x\\ 0\end{bmatrix}\\
 &\quad+\tilde\psi(\theta_{\mathrm{enc}},\mathbf y^{\mathrm n}_\tau,\mathbf u^{\mathrm n}_\tau),
\end{aligned}
\label{eq:aug_encoder}
\end{equation}
where $\mathbf y_\tau=\col(y_{\tau-n},\dots,y_\tau)\in\R^{3(n+1)}$ and
$\mathbf u_\tau=\col(u_{\mathrm{act},\tau-n},\dots,u_{\mathrm{act},\tau})\in\R^{3(n+1)}$
are the recorded histories over the encoder lag $n$,
$\mathbf y^{\mathrm n}_\tau$ and $\mathbf u^{\mathrm n}_\tau$ their normalised
counterparts, $\bar T_y=I_{n+1}\otimes T_y$, $\bar T_u=I_{n+1}\otimes T_u$,
$W^b\in\R^{6\times6(n+1)}$, $W^a\in\R^{n_{\bar x}\times6(n+1)}$,
$\tilde\psi:\R^{6(n+1)}\to\R^{6+n_{\bar x}}$ a multilayer perceptron, and
$\theta_{\mathrm{enc}}$ collects $W^b$, $W^a$ and the parameters of
$\tilde\psi$.
The encoder estimates all states, the measured positions included.
\todo{Missing: whether to mention the alternative of initialising the positions from the measurement~\cite[Remark~5.3]{kessels2025ai}.}
\todo{Align the figure's past-only encoder-history labels with
\eqref{eq:aug_encoder}, which also uses the current sample.}
\todo{VERIFY bib: \texttt{hoekstra2026encoder} = arXiv:2602.13108v1
(13 Feb 2026), Hoekstra, Gy\"or\"ok, T\'oth, Schoukens; matches PDF p.~1.
Check for a published (IFAC) version before submission.}

The physical rows start from the baseline, following the model-based
initialisation of Hoekstra et al.~\cite[Sec.~3.1]{hoekstra2026encoder}.
For zero input, every rest state with $\Theta=0$ is an equilibrium of
\eqref{eq:eom}, because $K$ in \eqref{eq:damping_stiffness_matrices} vanishes
on $X$ and $Y$.
We linearise the baseline with its nominal parameters at rest at $Y=0$.
Unlike~\cite[Eq.~(30)]{hoekstra2026lfr}, which fits the baseline encoder to
simulated baseline states, we compute $W^b$ from this linearisation.
The linearisation is known, so no encoder pretraining is needed.

Discretising the linearisation by zero-order hold over $T_s$ in scaled
coordinates gives $(A_d,B_d,C_d)$, whose reconstructability
map~\cite[Eqs.~(15) to (17)]{hoekstra2026encoder} is
\begin{equation}
 W^b=\begin{bmatrix}
  A_d^nO_n^{+} & \;\Gamma_n-A_d^nO_n^{+}H_n
 \end{bmatrix},
 \label{eq:aug_wb}
\end{equation}
where $O_n\in\R^{3(n+1)\times6}$ is the observability matrix of $(A_d,C_d)$
over $n+1$ samples, $H_n\in\R^{3(n+1)\times3(n+1)}$ the Toeplitz matrix of its
Markov parameters, $\Gamma_n\in\R^{6\times3(n+1)}$ the input-to-state map over the
same samples, and $(\cdot)^+$ the pseudoinverse.
The measured positions make $(A_d,C_d)$ observable, so $O_n$ has full column
rank although the baseline is marginally stable.

Only $W^b$ carries baseline information at the start.
$W^a$ is drawn by Xavier initialisation, as
in~\cite[Eq.~(31)]{hoekstra2026lfr}, and the output layer of $\tilde\psi$
starts at zero.
The encoder is trained jointly with the model, so the linearisation sets only
the starting value of $W^b$.
Such a model-based start gives better starting values and faster convergence
than a random encoder, at comparable final
accuracy~\cite[Sec.~4.4]{hoekstra2026encoder}.
The encoder paper shows this for a stable baseline whose parameters are not
co-estimated~\cite[Sec.~2, footnote~1]{hoekstra2026encoder}.
Here, $W^b$ uses the nominal parameters also when $\theta_{\mathrm{base}}$
starts detuned.
\todo{Missing: why linearise at $Y=0$. The encoder paper also linearises ``around the 0 point''~\cite[Sec.~4.2]{hoekstra2026encoder} but gives no reason.}
\todo{Missing: whether the nominal parameters in $W^b$ are acceptable prior information when $\theta_{\mathrm{base}}$ starts detuned, given that the black box of Section~\ref{sec:setup} gets no physical knowledge (D-243). Candidate: state it as a known inconsistency between encoder and physical block at the start (README sources.md trap~1).}

\subsection{Closed-loop identification}\label{sec:aug_closed_loop}
\ifdraft
\begin{itemize}
\item The records are measured under feedback and the model must predict the machine under feedback; the free axes make the open-loop gantry marginally stable, which violates the stability condition behind SUBNET consistency (Hoekstra 2026 Sec.~5.5, Beintema Cond.~1); the loop with the controller can satisfy it
\item We simulate each window in closed loop with the known controller, as Kessels does (Eqs.~(5.12), (5.13), Ch.~5; open-loop attempt failed, Remark~5.6); for linear systems the indirect approach (Forssell and Ljung); a motivation, not a guarantee
\item Both loops share $r$, $u^{\mathrm{ff}}$ and the controller; subtracting them gives our residual form; no plant feedthrough, so no algebraic loop
\item The residual form equals the machine loop only for a compatible controller and signals; the controller runs re-discretised at the model rate (difference to report (?))
\item $x^{c}_\tau=0$ needs no controller estimate and matches the encoder; Kessels' direct form needs the machine's controller state (Remark~5.4)
\item Training coordinates of $\theta_{\mathrm{base}}$ (moved from Section~II-C): nine log coordinates, bounded $m_\Delta$; every iterate satisfies \eqref{eq:lfr_admissibility}, so $M(Y)\succ0$ without clamp or penalty; positive-definite inertia, not positive raw masses (D-229)
\item Identification problem in one display (Hoekstra 2026 Eq.~(22) form): every window sample scored, closed loop, $\xi$, $\theta_{\mathrm{aug}}$, $\theta_{\mathrm{enc}}$ joint, controller fixed; unlike Kessels Eq.~(5.12), $\theta_{\mathrm{base}}$ is estimated
\item Adam then L-BFGS polish (Drenth Sec.~5)
\item Best validation free-run RMS selects; the free run carries pre-window error; horizons are separate quantities; polish kept if it improves and no meter regresses, the combination-error meter uses the true parameters (?)
\item A closed-loop fit can hide plant error (Results)
\item $\mathcal M_{\mathrm U}$: the model identified from the zero start (notation (?))
\end{itemize}
\fi

The records are measured under feedback, and the model must predict the
machine under feedback.
A rollout driven by the recorded input, as in the criterion
of~\cite[Eq.~(22)]{hoekstra2026lfr}, also meets a stability problem.
$X$ and $Y$ have no stiffness in \eqref{eq:damping_stiffness_matrices}, so the
open-loop gantry is marginally stable.
A velocity error then leaves a permanent position offset.
The consistency of the identification method~\cite[Sec.~5.5]{hoekstra2026lfr}
rests on that of SUBNET, which assumes an incrementally exponentially output
stable data-generating system~\cite[Cond.~1]{beintema2023subnet}.
The open-loop gantry does not satisfy this condition, whereas the loop with the
stabilising controller can.
\todo{Missing: evidence on the thesis data that open-loop rollouts accumulate position error on $X$ and $Y$; measure and report in Results.}

We therefore simulate each window in closed loop with the known controller.
Kessels trains augmented models in this way~\cite[Eqs.~(5.12),
(5.13)]{kessels2025ai}, which lets him update models of systems that are
unstable in open loop~\cite[Ch.~5]{kessels2025ai}.
He also reports that identifying the open-loop model from his closed-loop data
failed~\cite[Remark~5.6]{kessels2025ai}.
For linear systems, identifying the closed loop from the reference with a known
controller is the indirect approach~\cite{forssell1999revisited}.
Its output-error criterion remains consistent because the reference is
uncorrelated with the noise.
For our nonlinear model, these arguments motivate the choice but do not prove
consistency.

The controller of each record is known, linear and time invariant, with
state-space matrices $(A_{\mathrm{fb}},B_{\mathrm{fb}},C_{\mathrm{fb}},D_{\mathrm{fb}})$
that map the position error in actuator coordinates to actuator forces through
a controller state in $\R^{n_{\mathrm c}}$.
The machine and the model are driven by the same reference $r$ and feedforward
$u^{\mathrm{ff}}$ (Fig.~\ref{fig:aug_closed_loop}).
With $x^{\mathrm{fb}}$ the state of the machine's controller, the machine
applied the first two lines below, and the model applies the same controller to
$r_k-\hat y_k$.
Subtracting the machine's equations from the model's removes $r$ and
$u^{\mathrm{ff}}$ and gives our residual form in the last two lines:
\begin{equation}
\begin{aligned}
 x^{\mathrm{fb}}_{k+1}&=A_{\mathrm{fb}}x^{\mathrm{fb}}_k+B_{\mathrm{fb}}(r_k-y_k),\\
 u_{\mathrm{act},k}&=C_{\mathrm{fb}}x^{\mathrm{fb}}_k+D_{\mathrm{fb}}(r_k-y_k)+u^{\mathrm{ff}}_k,\\
 x^{c}_{k+1}&=A_{\mathrm{fb}}x^{c}_k+B_{\mathrm{fb}}(y_k-\hat y_k),\qquad x^{c}_\tau=0,\\
 \hat u_{\mathrm{act},k}&=u_{\mathrm{act},k}+C_{\mathrm{fb}}x^{c}_k+D_{\mathrm{fb}}(y_k-\hat y_k),
\end{aligned}
\label{eq:aug_residual_loop}
\end{equation}
where $r_k,u^{\mathrm{ff}}_k\in\R^3$,
$x^{c}=\hat x^{\mathrm{fb}}-x^{\mathrm{fb}}$ is the difference between the
controller states of the model and the machine, and
$\hat u_{\mathrm{act},k}\in\R^3$ is the model input, so that
$u_k=P\hat u_{\mathrm{act},k}$ in \eqref{eq:aug_transition}.
Since $h_{\mathrm{base}}$ has no feedthrough, the controller feedthrough
$D_{\mathrm{fb}}$ closes no algebraic loop.

\begin{figure}[t]
 \centering
 \includegraphics[width=\columnwidth]{closed_loop_same_reference.pdf}
 \caption{Recorded plant loop (top) and augmented-model loop (bottom),
 driven by one shared reference $r_k$ and feedforward $u_k^{\mathrm{ff}}$,
 with the same feedback controller. The model output is evaluated
 before the input advances its state to the next sample.}
 \label{fig:aug_closed_loop}
\end{figure}

The residual form equals the shared-reference loop when the recorded input was
produced by this controller from the same output samples.
The model loop runs the controller re-discretised at the model rate, whereas
the records were generated at a higher rate, so the two loops differ slightly.
\todo{Missing: where the measured difference between the two loops is reported (D-141: re-discretisation moves the peak output sensitivity at 150~Hz by 15.3\,\% and the phase margin by 3.6~degrees). Candidate: one sentence in Section~\ref{sec:setup} next to the controller rates.}
\todo{Check for Setup: run the re-discretised controller on the recorded $r-y$ and compare its feedback force with the recorded one. A rollout with $\hat y=y$ gives $\hat u=u$ by construction and checks nothing.}

Because $x^{c}$ is a difference of controller states, it needs no estimate.
Setting $x^{c}_\tau=0$ places the model's controller in the machine's
controller state at every window start, as the encoder places the model state
on the recorded trajectory.
Kessels' direct form instead needs the machine's controller state at every
window start, read from the machine or reconstructed from the reference, the
output and the controller~\cite[Remark~5.4]{kessels2025ai}.
Model error made before $\tau$ is therefore not carried into the window.

We estimate $\theta_{\mathrm{base}}$ in coordinates that keep
\eqref{eq:lfr_admissibility}, and with it well-posedness, at every iterate, as
Section~\ref{sec:params} requires.
They guarantee positive-definite inertia, not positive component masses.
The nine positive combinations are trained in logarithmic coordinates and the
signed $m_\Delta$ through a bounded map,
\begin{equation}
\begin{aligned}
 \theta_{\mathrm{base},j}&=\theta_{\mathrm{base},0,j}\,e^{\xi_j},
 \qquad j\neq8,\\
 m_\Delta&=\rho\,\frac{2}{L_b}\sqrt{m_\Sigma J_{\mathrm{eff}}}\,
   \tanh(\eta_0+s\,\xi_\Delta),
\end{aligned}
\label{eq:mdelta_map}
\end{equation}
where $\xi\in\R^{10}$ collects the training coordinates,
$\xi_\Delta=\xi_8$ belongs to $m_\Delta$, the eighth entry of \eqref{eq:phi},
$\theta_{\mathrm{base},0}$ is the start, $\rho\in(0,1)$ a margin, and
$\eta_0,s\in\R$ make $\xi=0$ return $\theta_{\mathrm{base},0}$ with
$\partial m_\Delta/\partial\xi_\Delta=m_{\Delta,0}$ there
(Appendix~\ref{app:lfr-wellposed}).
Since $\rho<1$ and $|\tanh|<1$, every $\xi$ satisfies
\eqref{eq:lfr_admissibility}, so $M(Y)\succ0$ for every $Y$ without a clamp
or penalty.

Over the window starts $\mathcal T$ and the window length $n_f$, the parameters
$\vartheta=(\xi,\theta_{\mathrm{aug}},\theta_{\mathrm{enc}})$ solve the
closed-loop form of the truncated criterion
of~\cite[Eq.~(22)]{hoekstra2026lfr},
\begin{equation}
\begin{aligned}
 &\min_{\vartheta}\;V(\vartheta)=\frac{1}{|\mathcal T|\,n_f\,n_y}
 \sum_{\tau\in\mathcal T}\sum_{\ell=0}^{n_f-1}
 \norm{y^{\mathrm n}_{\tau+\ell}-\hat y^{\mathrm n}_{\tau+\ell|\tau}}_2^2
 \quad\text{s.t.}\\
 &\tilde x_{k+1}=f_{\mathrm{base}}(\theta_{\mathrm{base}}(\xi),\tilde x_k,u_k)
   +f_{\mathrm{aug}}(\theta_{\mathrm{aug}},\hat x_k,u_k),\\
 &\bar x_{k+1}=g_{\mathrm{aug}}(\theta_{\mathrm{aug}},\hat x_k,u_k),\qquad
   \hat y_k=h_{\mathrm{base}}(\tilde x_k),\\
 &\col\bigl(T_x(\tilde x_{\tau}-\mu_x),\bar x_{\tau}\bigr)
   =\psi(\theta_{\mathrm{enc}},\mathbf y_\tau,\mathbf u_\tau),\\
 &x^{c}_{k+1}=A_{\mathrm{fb}}x^{c}_k+B_{\mathrm{fb}}(y_k-\hat y_k),\qquad x^{c}_\tau=0,\\
 &u_k=P\bigl(u_{\mathrm{act},k}+C_{\mathrm{fb}}x^{c}_k+D_{\mathrm{fb}}(y_k-\hat y_k)\bigr),
\end{aligned}
\label{eq:aug_loss}
\end{equation}
where $k$ runs over $\tau,\dots,\tau+n_f-1$ for every $\tau\in\mathcal T$,
$\theta_{\mathrm{base}}(\xi)$ is \eqref{eq:mdelta_map},
$\hat y_{\tau+\ell|\tau}$ is $\hat y_{\tau+\ell}$ of the window that starts at
$\tau$, the superscript $\mathrm n$ denotes normalisation as in
\eqref{eq:aug_norm}, and $n_y=3$.
Unlike Kessels, who keeps the first-principles parameters
fixed~\cite[Eq.~(5.12)]{kessels2025ai}, we estimate $\theta_{\mathrm{base}}$
jointly with the network and the encoder.

We minimise $V$ with Adam on minibatches of windows and then polish with
L-BFGS on a fixed set of windows, the order used by Drenth et
al.~\cite[Sec.~5]{drenth2025lpvlfr}.

The selected checkpoint has the lowest free-run error on the validation
records, the record-length-weighted mean of the per-record RMS error in metres.
The free run simulates \eqref{eq:aug_loss} over each whole record, with the
encoder and $x^{c}=0$ applied once at its start.
Model error therefore accumulates over the record.
The encoder lag $n$, the window length $n_f$ and this free-run horizon are
separate quantities.
The polished parameters replace the selected checkpoint only if they lower the
same error and no monitored quantity of the polish worsens by more than a
relative tolerance.
One of these quantities is the combination error, the root mean square of the
relative errors of the ten combinations with respect to their true values, so
the rule uses the true parameters.
\todo{Missing: an acceptance rule without the true parameters. The polish meter \texttt{combo\_err} (\texttt{GD/training.py::\_gantry\_meters}, D-171) compares against the \texttt{gantry\_ss} truth, an oracle inside training (sources.md trap~6). Candidate: drop the meter and keep only the validation condition, or report the rule as an oracle choice in Section~\ref{sec:setup}.}

A good closed-loop fit does not by itself show that the model is correct in
open loop, because the controller acts on the output error and can mask plant
error.
Section~\ref{sec:results} therefore also evaluates the model outside the
training loop.

We call the model \eqref{eq:aug_transition} identified by \eqref{eq:aug_loss}
from the zero start \eqref{eq:aug_zero_init}, the encoder start of
Section~\ref{sec:aug_encoder} and $\theta_{\mathrm{base},0}$, and selected as
above, $\mathcal M_{\mathrm U}$.
Its inputs are the recorded actuator forces and positions, and its output is
the predicted actuator positions $\hat y$.
\todo{Missing: model notation. Candidate: $\mathcal M_{\mathrm U}$ and $\mathcal M_{\mathrm{PS2}}$, since $M$ is the inertia matrix (README open conflict on model names).}

\subsection{Added-state initialisation}\label{sec:aug_ps2}
\ifdraft
\begin{itemize}
\item At the zero start the added states do not reach the output, so $V$ gives their encoder rows and the added-state rows of $\phi_{\mathrm{aug}}$ no gradient; plain training then learns static corrections or real lags (Results) (PS2-METHOD Sec.~2.2)
\item PS2 (this work) adapts two-stage regression (Hefny Sec.~3) and its use before BPTT (Downey Sec.~4.2); unlike Hefny's first stage, the target is the current model's residual
\item S1: decoder and the encoder's added-state rows fitted to the current closed-loop residual over the encoder horizon (\texttt{PS2Opening.\_s1\_step})
\item S2: $g_{\mathrm{aug}}$ fitted once to propagate the encoder states one sample (\texttt{PS2Opening.\_fit\_map})
\item S3: decoder discarded, Section~III-C continues unchanged; the screen skips S2 when the residual is not predictable; plateau or cap ends S1 and S2 (\texttt{PS2Spec}); values in Section~V
\item $\mathcal M_{\mathrm{PS2}}$: $\mathcal M_{\mathrm U}$ initialised by PS2 (notation (?))
\end{itemize}
\fi

At the zero start \eqref{eq:aug_zero_init}, $\hat y$ does not depend on the
added states, so $V$ gives $W^a$ and the added-state rows of
$\phi_{\mathrm{aug}}$ no gradient.
Plain training afterwards learns static corrections or real lags instead of the
missing dynamics (Section~\ref{sec:results}).

We propose PS2, which adapts the two-stage regression of Hefny et
al.~\cite[Sec.~3]{hefny2015supervised} and its use as an initialisation before
backpropagation through time~\cite[Sec.~4.2]{downey2017psrnn}.
Unlike their first stage, which regresses features of future observations,
stage S1 fits a temporary decoder and the encoder's added-state rows to the
residual of the current model over the next $n+1$ samples:
\begin{equation}
\begin{aligned}
 &\min_{\theta_\lambda,\,\theta^{a}_{\mathrm{enc}}}\;
 \sum_{\tau\in\mathcal T}
 \norm{\lambda(\theta_\lambda,\bar x_{\tau|\tau})-\varepsilon_\tau}_2^2,\\
 &\varepsilon_\tau=\kappa\,\col\bigl(y^{\mathrm n}_{\tau+\ell}-\hat y^{\mathrm n}_{\tau+\ell|\tau}\bigr)_{\ell=0}^{n},
\end{aligned}
\label{eq:ps2_s1}
\end{equation}
where $\lambda:\R^{n_{\bar x}}\to\R^{3(n+1)}$ is a perceptron with one tanh
hidden layer and parameters $\theta_\lambda$, $\theta^{a}_{\mathrm{enc}}$ are
the encoder parameters that enter only $\bar x_{\tau|\tau}$ of
\eqref{eq:aug_encoder}, $\hat y_{\tau+\ell|\tau}$ is the closed-loop prediction
of \eqref{eq:aug_loss} at the current parameters, and $\kappa>0$ scales
$\varepsilon_\tau$ to unit root mean square over the windows.
S1 thus trains $\bar x_{\tau|\tau}$ to predict, from the history, what the
current model misses.

Stage S2 then fits the deployed added-state transition, once, to propagate
these states one sample ahead:
\begin{equation}
 \min_{\theta_{\mathrm{aug}}}\;\frac{1}{\bar\sigma^2}\sum_{\tau\in\mathcal T}
 \norm{g_{\mathrm{aug}}(\theta_{\mathrm{aug}},\hat x_{\tau|\tau},u_\tau)
 -\bar x_{\tau+1|\tau+1}}_2^2,
 \label{eq:ps2_s2}
\end{equation}
where $\hat x_{\tau|\tau}$ and $\bar x_{\tau+1|\tau+1}$ are encoder outputs
\eqref{eq:aug_encoder} at consecutive window starts with $\theta_{\mathrm{enc}}$
fixed, $u_\tau=Pu_{\mathrm{act},\tau}$ is the recorded input, and
$\bar\sigma^2$ is the mean variance of the entries of
$\bar x_{\tau+1|\tau+1}$.
S2 thus gives $g_{\mathrm{aug}}$ the dynamics of the S1 states by regression,
without the gradient through $V$ that the zero start blocks.

In S3, the decoder is discarded and the identification of
Section~\ref{sec:aug_closed_loop} continues from the result, with the model,
the criterion $V$, the optimiser and the selection unchanged.
S2 is skipped, and S3 starts directly, if the fraction $\gamma$ of the residual
energy that $\lambda$ leaves unexplained, on windows of training records
withheld from S1 and S2, exceeds a threshold $\bar\gamma$ at a fixed update or
at the end of S1.
After a minimum number of updates, S1 ends when $\gamma$ improves by less than a
relative tolerance over successive checks, or at an update cap, and S2 by the
same rule on its normalised residual on the withheld windows.

We call $\mathcal M_{\mathrm U}$ initialised by PS2 $\mathcal M_{\mathrm{PS2}}$,
with the inputs, output, criterion and selection of $\mathcal M_{\mathrm U}$.
